\chapter{Optimal Transport のエントロピー正則化}
\label{c-entropic}

本章では、Kantorovich 定式化の optimal transport およびその多くの一般化に対する近似解を求めるための数値スキームの一族を導入する。これはもとの問題にエントロピー正則化ペナルティを加えることで実現される。この正則化にはいくつかの重要な利点があり、それらを総合すると非常に有用なツールとなる：正則化された問題の最小化は単純な交互最小化スキームで解くことができ、そのスキームは単純な行列-ベクトル積の反復に帰着するため、GPU 上での実行に特に適している。また、ある応用では $n\times m$ コスト行列を格納する必要がなく、カーネル評価へのアクセスのみで済む。さらに、大規模な測度群が同じ台を共有する場合、これらの行列-ベクトル積はすべて行列-行列積として実行でき、大幅な高速化が得られる。結果として得られる近似距離は入力ヒストグラムの重みおよび Dirac の位置に関して滑らかであり、自動微分を用いて微分することができる。

%%%%%%%%%%%%%%%%%%%%%%%%%%%%%%%%%%%%%%%%%%%
\section{エントロピー正則化}

カップリング行列の離散エントロピーは以下で定義される。
\eql{\label{eq-discr-entropy}
	\HD(\P) \eqdef -\sum_{i,j} \P_{i,j} (\log(\P_{i,j})-1),
}
ベクトルに対しても同様に定義され、$\a_j$ のいずれかが $0$ または負の場合は $\HD(\a) = -\infty$ とする規約を用いる。関数 $\HD$ は $1$-強凹関数である。これはその Hesse 行列が $\partial^2 \HD(P)=-\diag(1/\P_{i,j})$ であり、$\P_{i,j} \leq 1$ であるためである。Optimal transport のエントロピー正則化の考え方は、もとの輸送問題~\eqref{eq-mk-discr}の近似解を得るために $-\HD$ を正則化関数として用いることである：
\eql{\label{eq-regularized-discr}
	\MKD_\C^\varepsilon(\a,\b) \eqdef
	\umin{\P \in \CouplingsD(\a,\b)}
		\dotp{\P}{\C} - \varepsilon \HD(\P).
}
目的関数は $\varepsilon$-強凸関数であるため、問題 (\ref{eq-regularized-discr}) は一意の最適解を持つ。Optimal transport 問題をエントロピー項で正則化する考え方は、輸送理論におけるモデリングの着想~\citep{wilson1969use}に遡ることができる：ネットワークにおける実際の交通パターンは、optimal transport 問題の解によって予測されるものとは一致しない。実際、~\eqref{eq-mk-discr}に対する最適カップリングのスパース性の結果として、後者が少数の経路に依存する傾向があるのに対し、前者はより拡散的である。このスパース性を緩和するため、輸送分野の研究者たちは、周辺分布と輸送コストが与えられた場合に、より「ぼやけた」交通予測を形成できる「重力」モデル~\citep{erlander1980optimal}と呼ばれるモデルを提案した。

図~\ref{fig-impact-eps}は、単体 $\simplex_3$（したがって2次元の三角形として可視化できる）上の線形計画に対してエントロピーが正則化する効果を示している。エントロピーがもとの LP 解を三角形の境界から押し離す様子に注目されたい。最適な $\P_\varepsilon$ は漸進的に三角形の「エントロピー中心」に向かって移動する。これは以下の命題でさらに詳述される。正則化された問題の解がもとの線形計画の最適解に収束することは \citet{CominettiAsympt} によって研究されており、精密な漸近挙動が与えられている。


\begin{figure}[h!]
\centering
\includegraphics[width=.7\linewidth]{entropic/simplex}
\caption{\label{fig-impact-eps}
$\P_\varepsilon = \argmin_{\P \in \simplex_3} \dotp{\C}{\P}-\varepsilon\HD(\P)$ を $\varepsilon$ を変化させて解いた場合の、単体上の線形関数の最適化に対する $\varepsilon$ の影響。
}
\end{figure}


\begin{prop}[$\varepsilon$ に関する収束]\label{prop-convergence-eps}
\eqref{eq-regularized-discr}の一意の解 $\P_\varepsilon$ は、Kantorovich 問題のすべての最適解の集合の中でエントロピーが最大の最適解に収束する。すなわち
\eql{\label{eq-entropy-conv-1}
	\P_\varepsilon \overset{\varepsilon \rightarrow 0}{\longrightarrow}
	\uargmin{\P} \enscond{ -\HD(\P) }{
		\P \in \CouplingsD(\a,\b), \dotp{\P}{\C} = \MKD_\C(\a,\b),
	}
}
したがって特に
\eq{
	\MKD_\C^\varepsilon(\a,\b) \overset{\varepsilon \rightarrow 0}{\longrightarrow} \MKD_\C(\a,\b).
}
さらに以下が成り立つ。
\eql{\label{eq-entropy-conv-2}
	\P_\varepsilon \overset{\varepsilon \rightarrow \infty}{\longrightarrow}
	\a \otimes \b = \a \transp{\b} = (\a_i \b_j)_{i,j}.
}
\end{prop}
\begin{proof}
	 $\varepsilon_\ell \rightarrow 0$ かつ $\varepsilon_\ell > 0$ であるような列 $(\varepsilon_\ell)_\ell$ を考える。
 	$\P_\ell$ を $\varepsilon=\varepsilon_\ell$ に対する~\eqref{eq-regularized-discr}の解とする。
	$\CouplingsD(\a,\b)$ は有界であるから、$\P_\ell \rightarrow \P^\star$ となる部分列を抽出できる（簡潔さのため部分列にラベルを付け替えない）。$\CouplingsD(\a,\b)$ は閉であるから、$\P^\star \in \CouplingsD(\a,\b)$ である。$\dotp{\C}{\P} = \MKD_\C(\a,\b)$ を満たす任意の $\P$ を考える。$\P$ と $\P_\ell$ がそれぞれの最適化問題（$\varepsilon=0$ と $\varepsilon=\varepsilon_\ell$）の最適解であることから、
 	\eql{\label{eq-proof-gamma-conv}
 		0 \leq \dotp{\C}{\P_\ell} - \dotp{\C}{\P} \leq \varepsilon_\ell ( \HD(\P_\ell)-\HD(\P) ).
 	}
 	$\HD$ は連続であるから、この式で $\ell \rightarrow +\infty$ とすると、$\dotp{\C}{\P^\star} = \dotp{\C}{\P}$ となり、$\P^\star$ は~\eqref{eq-entropy-conv-1}の実行可能点である。さらに、\eqref{eq-proof-gamma-conv}を $\varepsilon_\ell$ で割って極限をとると、$\HD(\P) \leq \HD(\P^\star)$ が得られ、$\P^\star$ が~\eqref{eq-entropy-conv-1}の解であることがわかる。$-\HD$ の狭義凸性によりこの問題の解 $\P_0^\star$ は一意であるから、$\P^\star = \P_0^\star$ であり、列全体が収束する。
	$\varepsilon\rightarrow+\infty$ の極限では、同様の証明により、むしろ以下の問題を考えるべきであることが示される。
	\eq{
		\umin{\P \in \CouplingsD(\a,\b)} -\HD(\P),
	}
	この問題の解は $\a \otimes \b$ である。
\end{proof}

式~\eqref{eq-entropy-conv-1}は、小さな正則化 $\varepsilon$ に対して、解が最大エントロピーの最適輸送カップリングに収束することを述べている。
対照的に、\eqref{eq-entropy-conv-2}は、大きな正則化に対して、解が所定の周辺分布 $\a,\b$ の間の最大エントロピーのカップリング、すなわち $\a,\b$ に従って分布する2つの独立な確率変数間の同時確率に収束することを示している。
この収束の精密な解析は \citet{CominettiAsympt} によって行われており、$\varepsilon=0$（および $\varepsilon=+\infty$）近傍における $\varepsilon$（および $1/\varepsilon$）の1次展開を含む。
図~\ref{fig-entropic-densities}および図~\ref{fig-entropic}は、これら2つの収束の効果を視覚的に示している。重要な洞察は、$\varepsilon$ が増加するにつれて、最適カップリングのスパース性が低下し（所定の閾値より大きい成分を持つという意味で）、それが計算アルゴリズムの高速化（第~\ref{sec-sinkhorn}節で研究する）と統計的収束の加速（第~\ref{sec-entropy-ot-mmd}節で示す）の両方をもたらすということである。

\newcommand{\myfigSinkEps}[1]{%
\imgBox{\includegraphics[width=.23\linewidth]{sinkhorn-eps/evol-img-#1}}}
\newcommand{\myfigSinkEpsV}[1]{\includegraphics[width=.24\linewidth]{sinkhorn-eps/evol-3d-#1}}

\begin{figure}[h!]
\centering
\begin{tabular}{@{}c@{}c@{}c@{}c@{}}
\myfigSinkEps{1} &
\myfigSinkEps{3} &
\myfigSinkEps{4} &
\myfigSinkEps{5} \\
\myfigSinkEpsV{1} &
\myfigSinkEpsV{3} &
\myfigSinkEpsV{4} &
\myfigSinkEpsV{5} \\
$\varepsilon=10$ &
$\varepsilon=1$ &
$\varepsilon=10^{-1}$ &
$\varepsilon=10^{-2}$
\end{tabular}
\caption{\label{fig-entropic-densities}
命題~\ref{prop-convergence-eps}を例示する、2つの1次元密度間のカップリングに対する $\varepsilon$ の影響。
上段：2つの1次元密度間。下段：同数 $n=m$ の点を持つ2つの2次元離散経験密度間（最適な $(\P_{i,j})_{i,j}$ のうち小さな閾値を超える成分のみが $x_i$ と $y_j$ の間の線分として表示されている）。
}
\end{figure}


\begin{figure}[h!]
\centering
\includegraphics[width=\linewidth]{entropic/matching-2d}
\caption{\label{fig-entropic}
同数 $n=m$ の点を持つ2つの2次元離散経験密度間のカップリングに対する $\varepsilon$ の影響（最適な $(\P_{i,j})_{i,j}$ のうち小さな閾値を超える成分のみが $x_i$ と $y_j$ の間の線分として表示されている）。
}
\end{figure}

カップリング間の Kullback--Leibler ダイバージェンスを以下で定義する。
\eql{\label{eq-kl-defn}
	\KLD(\P|\K) \eqdef \sum_{i,j}  \P_{i,j} \log\pa{\frac{\P_{i,j}}{\K_{i,j}}} - \P_{i,j} + \K_{i,j},
}
\eqref{eq-regularized-discr}の一意の解 $\P_\varepsilon$ は、コスト行列 $\C$ に付随する Gibbs カーネルの $\CouplingsD(\a,\b)$ への射影であり、
\eq{
	\K_{i,j} \eqdef e^{-\frac{\C_{i,j}}{\varepsilon}}.
}
実際、上の定義を用いると以下が成り立つ。
\eql{\label{eq-kl-proj}
	\P_\varepsilon = \Proj_{\CouplingsD(\a,\b)}^\KLD(\K) \eqdef \uargmin{\P \in \CouplingsD(\a,\b)} \KLD(\P|\K).
}




%%%%%
\begin{rem1}{離散測度間のエントロピー正則化}\label{rem:entrop-reg}
\eqref{eq-discr-meas}の形の離散測度に対して、正則化輸送の定義は自然に以下のように拡張される。
\eql{\label{eq-entropic-disc-meas}
	\MK_\c^\varepsilon(\al,\be) \eqdef \MKD_\C^\varepsilon(\a,\b),
}
ここで $\C_{i,j}=\c(x_i,y_j)$ であり、入力測度を支える位置 $(x_i,y_j)$ への依存性を強調している。
\end{rem1}
%%%%%


%%%%%
\begin{rem2}{一般的な定式化}
任意の測度を考えるには、離散エントロピーを積測度 $\d\al\otimes\d\be(x,y) \eqdef \d\al(x)\d\be(y)$ に関する相対エントロピーに置き換え、~\eqref{eq-mk-generic}の正則化版を以下のように提案する。
\eql{\label{eq-entropic-generic}
	\MK_\c^\varepsilon(\al,\be) \eqdef
	\umin{\pi \in \Couplings(\al,\be)}
		\int_{\X \times \Y} c(x,y) \d\pi(x,y) + \varepsilon \KL(\pi|\al\otimes\be),
}
ここで相対エントロピーは離散 Kullback--Leibler ダイバージェンス~\eqref{eq-kl-defn}の一般化であり、
\eql{\label{eq-defn-rel-entropy}
\begin{aligned}	 \KL(\pi|\xi) \eqdef \int_{\X \times \Y} \log\Big( \frac{\d \pi}{\d\xi}(x,y) \Big) \d\pi(x,y)\\
	  + \int_{\X \times \Y} (\d\xi(x,y)-\d\pi(x,y)),
	  \end{aligned}
}
$\pi$ が $\xi$ に関する密度 $\frac{\d \pi}{\d\xi}$ を持たない場合は $\KL(\pi|\xi)=+\infty$ と規約する。
~\eqref{eq-entropic-generic}でエントロピー正則化項 $\KL(\cdot|\al\otimes\be)$ を定義するために選ばれた参照測度 $\al\otimes\be$ は特別な役割を果たさず、その台のみが重要であることを認識しておくことが重要である。これは以下の命題によって示される。

\begin{prop}
	任意の $\pi \in \Couplings(\al,\be)$ に対し、$(\al,\be)$ と同じ零測度集合を持つ（したがって互いに密度を持つ）任意の $(\al',\be')$ に対して以下が成り立つ。
	\eq{
		\KL(\pi|\al\otimes\be) = \KL(\pi|\al'\otimes\be') - \KL(\al\otimes\be|\al'\otimes\be').
	}
\end{prop}

この命題は、\eqref{eq-entropic-generic}において $\KL(\cdot|\al\otimes\be)$ の代わりに $\KL(\cdot|\al'\otimes\be')$ を用いても同じ解が得られることを示している。


式~\eqref{eq-entropic-generic}は射影問題として再定式化できる。
\eql{\label{eq-entropic-generic-proj}
	\umin{\pi \in \Couplings(\al,\be)} \KL(\pi|\Kk),
}
ここで $\Kk$ は Gibbs 分布 $\d\Kk(x,y) \eqdef e^{-\frac{c(x,y)}{\varepsilon}} \d\mu(x)\d\nu(y)$ である。
この問題は「静的 Schr\"odinger 問題」~\citep{LeonardSchroedinger,RuschendorfThomsen}としばしば呼ばれる。これは Schr\"odinger が統計物理学において最初に考えた問題~\citep{Schroedinger31}に由来する。
$\varepsilon \rightarrow 0$ のとき、\eqref{eq-entropic-generic-proj}の一意の解は~\eqref{eq-mk-generic}の最大エントロピー解に収束する。\citep{leonard2012schrodinger,2017-carlier-SIMA}を参照されたい。
第~\ref{sec-entropic-dynamic}節では、2つの測度の点を結ぶ経路の空間上における Schr\"odinger 問題の別の「動的」定式化を詳述する。
\end{rem2}
%%%%%

%%%%%
\begin{rem2}{相互エントロピー}
\eqref{eq-ot-proba-interpretation}と同様に、確率変数を用いて~\eqref{eq-entropic-generic}を以下のように言い換えることができる。
\eq{
	\MK_\c^\varepsilon(\al,\be) = \umin{(X,Y)} \enscond{ \EE_{(X,Y)}(c(X,Y)) + \epsilon \text{I}(X,Y) }{ X \sim \al, Y \sim \be },
}
ここで $(X,Y)$ の分布を $\pi$ と書くと、$\text{I}(X,Y) \eqdef \KL(\pi|\al\otimes\be)$ は2つの確率変数間のいわゆる相互情報量である。$\text{I}(X,Y) \geq 0$ が成り立ち、$\text{I}(X,Y)=0$ であることと2つの確率変数が独立であることは同値である。
\end{rem2}

%%%%%
\begin{rem2}{独立性とカップリング}
	カップリング $\pi \in \Couplings(\al,\be)$ は、$(\X,\Y)$ 上で定義された確率変数の組 $(X,Y)$ の分布を記述する。ここで $X$（および $Y$）の法則はそれぞれ $\al$（および $\be$）である。
	命題~\ref{prop-convergence-eps}は一般的な（必ずしも離散でない）測度に対しても成り立つ。したがって、\eqref{eq-entropic-generic}の解 $\pi_\varepsilon$ は $\varepsilon \rightarrow+\infty$ のとき直積カップリング $\al \otimes \be$ に収束する。このカップリング $\al \otimes \be$ は確率変数 $(X,Y)$ が独立であることに対応する。
	対照的に、$\varepsilon \rightarrow 0$ のとき、$\pi_\varepsilon$ は OT 問題~\eqref{eq-mk-generic}の解 $\pi_0$ に収束する。$\X=\Y=\RR^\dim$ 上で、$\al$ と $\be$ が Lebesgue 測度に関する密度を持つ場合、注意~\ref{rem-exist-mongemap}で詳述されるように、$\pi_0$ は一意であり全単射 Monge 写像 $\T : \RR^\dim \rightarrow \RR^\dim$ のグラフ上に台を持つ。この場合、$Y=\T(X)$ かつ $X=\T^{-1}(Y)$ であるから、$(X,Y)$ はある意味で完全に従属的である。
	1次元の単純な場合 $\dim=1$ では、$X$ と $Y$ の間の従属構造を可視化する便利な方法は、同時分布 $\pi$ に付随するコピュラ $\xi_\pi$ を用いることである。\eqref{eq-cumul-defn}で定義された累積関数はカップリングに対して以下のように拡張される。
	\eq{
		\foralls (x,y) \in \RR^2, \quad
		\cumul{\pi}(x,y) \eqdef \int_{-\infty}^x \int_{-\infty}^y \d\pi.
	}
	コピュラは以下で定義される。
	\eq{
		\foralls (s,t) \in [0,1]^2, \quad
		\xi_{\pi}(s,t) \eqdef \cumul{\pi}( \cumul{\al}^{-1}(s),\cumul{\be}^{-1}(t) ),
	}
	ここで累積関数の疑似逆は~\eqref{eq-pseudo-inv-cum}で定義されている。
	独立な変数の場合、$\varepsilon=+\infty$、\ie $\pi=\al\otimes \be$ であり、$\xi_{\pi_{+\infty}}(s,t)=st$ となる。
	対照的に、完全に従属な変数の場合、$\varepsilon=+\infty$ であり、$\xi_{\pi_0}(s,t)=\min(s,t)$ となる。
	図~\ref{fig-entropic-copula}は、エントロピー正則化がこれら2つの極端な場合の間を補間するコピュラ $\xi_{\pi_\varepsilon}$ を生成する様子を示している。
\end{rem2}
%%%%%

\newcommand{\myfigSinkCopul}[1]{\imgBox{\includegraphics[width=.185\linewidth,trim=65 55 45 40,clip]{sinkhorn-copula/evol-#1}}}
\begin{figure}[h!]
\centering
\begin{tabular}{@{}c@{\hspace{5mm}}c@{}}
\includegraphics[width=.4\linewidth]{sinkhorn-copula/input-1} &
\includegraphics[width=.4\linewidth]{sinkhorn-copula/input-2}\\
{\color{red}$\al$} & {\color{blue}$\be$}
\end{tabular}
%
\begin{tabular}{@{}c@{}c@{}c@{}c@{}c@{}}
\myfigSinkCopul{levelsets-1} &
\myfigSinkCopul{levelsets-2} &
\myfigSinkCopul{levelsets-3} &
\myfigSinkCopul{levelsets-4} &
\myfigSinkCopul{levelsets-5} \\
\myfigSinkCopul{copula-1} &
\myfigSinkCopul{copula-2} &
\myfigSinkCopul{copula-3} &
\myfigSinkCopul{copula-4} &
\myfigSinkCopul{copula-5} \\
{\color[rgb]{1,0,1} $\varepsilon=10$} &
{\color[rgb]{.75,.25,.75} $\varepsilon=1$} &
{\color[rgb]{.5,.5,.5} $\varepsilon=0.5 \cdot 10^{-1}$} &
{\color[rgb]{.25,.75,.25} $\varepsilon=10^{-1}$} &
{\color[rgb]{0,1,0} $\varepsilon=10^{-3}$}
\end{tabular}
\caption{\label{fig-entropic-copula}
上段：~\eqref{eq-entropic-generic}の解 $\pi_\varepsilon$ の $\varepsilon$ に伴う変化。
下段：コピュラ関数 $\xi_{\pi_\varepsilon}$ の変化。
}
\end{figure}


%%%%%%%%%%%%%%%%%%%%%%%%%%%%%%%%%%%%%%%%%%%
\section{Sinkhorn アルゴリズムとその収束}
\label{sec-sinkhorn}

以下の命題は、\eqref{eq-regularized-discr}の解が $n+m$ 個の変数でパラメータ化できる特別な形を持つことを示す。このパラメータ化は本質的に双対的である。なぜなら、$\CouplingsD(\a,\b)$ に属するカップリング $\P$ は $nm$ 個の変数を持つが $n+m$ 個の制約を持つからである。

\begin{prop}\label{prop-regularized-primal}
\eqref{eq-regularized-discr}の解は一意であり、以下の形を持つ。
\eql{\label{eq-scaling-form}
	\foralls (i,j) \in \range{n} \times \range{m}, \quad \P_{i,j} = \uD_i \K_{i,j} \vD_j
}
ここで $(\uD,\vD) \in \RR_+^n \times \RR_+^m$ は2つの（未知の）スケーリング変数である。
\end{prop}

\begin{proof}
各周辺制約に対する2つの双対変数 $\fD\in\RR^n,\gD\in\RR^m$ を導入すると、\eqref{eq-regularized-discr}の Lagrangian は以下のようになる。
\eq{\label{eq-sinkhorn-lagrangian}
	\Lag(\P,\fD,\gD)= \dotp{\P}{\C} - \varepsilon \HD(\P) - \dotp{\fD}{\P\ones_m-\a}-\dotp{\gD}{\transp{\P}\ones_n-\b}.
}
1次の最適性条件から
$$
	\frac{\partial\Lag(\P,\fD,\gD)}{\partial \P_{i,j}}= \C_{i,j} + \varepsilon \log(\P_{i,j}) - \fD_i -\gD_j=0,
$$
が得られ、正則化問題の最適カップリング $\P$ に対して $\P_{i,j}=e^{\fD_i/\varepsilon}e^{-\C_{i,j}/\varepsilon}e^{\gD_j/\varepsilon}$ という表現が得られる。これは非負ベクトル $\uD$ と $\vD$ を用いて上記の形に書き直すことができる。
\end{proof}

%%%
\paragraph{行列スケーリングとしての正則化 OT}

式~\eqref{eq-scaling-form}で示された最適解の分解は、行列形式で $\P=\diag(\uD)\K\diag(\vD)$ と便利に書き直すことができる。
変数 $(\uD,\vD)$ は、$\CouplingsD(\a,\b)$ に内在する質量保存制約に対応する以下の非線形方程式を満たさなければならない：
\eql{\label{eq-dualsinkhorn-constraints}
	\diag(\uD)\K\diag(\vD)\ones_m=\a,
	\qandq
	\diag(\vD)\K^\top \diag(\uD)\ones_n=\b.
}
$\diag(\vD)\ones_m$ は単に $\vD$ であり、$\diag(\uD)$ と $\K \vD$ の積は以下のように簡略化される。
\eql{\label{eq-dualsinkhorn-constraints2}
	\uD \odot (\K \vD) = \a
	\qandq
	\vD \odot (\transp{\K}\uD) = \b,
}
ここで $\odot$ はベクトルの要素ごとの積に対応する。この問題は数値解析のコミュニティでは行列スケーリング問題として知られている（\citep{nemirovski1999complexity}およびそこでの参考文献を参照）。
これらの方程式を扱う直観的な方法は、まず式~\eqref{eq-dualsinkhorn-constraints2}の左辺を満たすように $\uD$ を修正し、次にその右辺を満たすように $\vD$ を修正することで、反復的に解くことである。これら2つの更新が Sinkhorn アルゴリズムを定義する。
\eql{\label{eq-sinkhorn}
	\itt{\uD} \eqdef \frac{\a}{\K \it{\vD}}
	\qandq
	\itt{\vD} \eqdef \frac{\b}{\transp{\K}\itt{\uD}},
}
初期値は任意の正ベクトル $\init{\vD} = \ones_m$ である。上で用いた2つのベクトル間の除算は要素ごとに理解されるべきである。異なる初期化は $\uD,\vD$ の異なる解につながりうる。これは $\uD,\vD$ が乗法的定数を除いてのみ定義されるためである（$\uD,\vD$ が~\eqref{eq-dualsinkhorn-constraints}を満たすならば、任意の $\lambda>0$ に対して $\lambda\uD,\vD/\lambda$ も満たす）。
しかしながら、これらの反復は収束し（反復射影による正当化については注意~\ref{rem-iterative-projection}を、厳密な縮小結果については注意~\ref{rem-global-conv-sinkh}を参照）、すべて同一の最適カップリング $\diag(\uD)\K\diag(\vD)$ に帰結する。
図~\ref{fig-sinkhorn-convergence}の上段は、Sinkhorn 反復により計算されるカップリング $\diag(\it{\uD})\K\diag(\it{\vD})$ の変化を示している。これは Gibbs カーネル $\K$ から出発し、質量を対角線から漸進的に移動させながら~\eqref{eq-regularized-discr}を解く最適カップリングに向かって変化する。

\begin{rem}[歴史的観点] 反復~\eqref{eq-sinkhorn}は~\citep{yule1912methods,kruithof}に最初に現れた。後に反復比例フィッティング手順（IPFP）\citet{DemingStephanIPFP}および RAS~\citep{bacharach1965estimating}法~\citep{ReviewSinkhorn}として知られるようになった。その収束の証明は \citet{Sinkhorn64} に帰せられ、このアルゴリズムの名称の由来となった。このアルゴリズムは後に \citet{Ruschendorf95} により無限次元に拡張された。この正則化は、重力モデル~\citep{wilson1969use,erlander1980optimal,erlander1990gravity}の名の下に、optimal transport 問題の近似解を得るために経済学の分野で用いられた。
割当の場合、すなわち $\a=\b=\ones_n/n$ の場合に、\citet{kosowsky1994invisible}により「softassign」と名づけられ、より最近では \citet{Galichon-Entropic} により経済学におけるマッチング問題を解くために用いられた。
この正則化は \citep{CuturiSinkhorn} 以降、データサイエンス（機械学習、コンピュータビジョン、グラフィックス、イメージングを含む）において改めて注目されている。Sinkhorn アルゴリズムが optimal transport の効率的かつスケーラブルな近似を提供すること、複数の OT 問題を同時に解く際のシームレスな並列化（特に GPU 上。注意~\ref{rem-parallel}を参照）、およびこの正則化された量が線形計画定式化とは異なり微分可能な損失関数を定義すること（第~\ref{sec-regularized-cost}節を参照）が示された。
Sinkhorn アルゴリズムの拡張や一般化は無数に存在する（例えば第~\ref{sec-generalized}節を参照）。例えば、$\a=\b$ の場合、対称性を保つために平均化された射影反復を用いることができる~\citep{knight2014symmetry}。
\end{rem}

\begin{figure}[h!]
\centering
\includegraphics[width=\linewidth]{entropic/sinkhorn-convergence}
\begin{tabular}{c}
\includegraphics[width=.5\linewidth]{entropic/sinkhorn-rates}\\
{\color{blue}$\varepsilon=10$} \quad {\color[rgb]{0.5,0,0.5}$\varepsilon=0.1$} \quad {\color{red}$\varepsilon=10^{-3}$}
\end{tabular}
\caption{\label{fig-sinkhorn-convergence}
上段：$\X=[0,1]$ 上の1次元密度、$c(x,y)=|x-y|^2$、$\varepsilon=0.1$ に対する Sinkhorn 反復の第 $\ell$ 反復で計算されたカップリング $\it{\pi_\varepsilon}=\diag(\it{\uD})\K\diag(\it{\vD})$ の変化。
下段：周辺制約違反 $\log( \|\it{\pi_\varepsilon} \ones_m - \b \|_1 )$ で測定した Sinkhorn の収束速度に対する $\varepsilon$ の影響。
}
\end{figure}


%%%%%%%%%%
\begin{rem1}{全体の計算量}\label{rem-complexity-rounding}
	慎重な収束解析（簡潔さのため $n=m$ を仮定）により、\citet{altschuler2017near}は $\varepsilon = \frac{4\log(n)}{\tau}$ と設定すると、$O(\norm{\C}_\infty^3 \log(n) \tau^{-3})$ 回の Sinkhorn 反復（有効なカップリング $\hat \P \in \CouplingsD(\a,\b)$ を計算するための追加の丸めステップを含む）で $\dotp{\hat \P}{\C} \leq \MKD_{\C}(\a,\b) + \tau$ を保証するのに十分であることを示した。これは Sinkhorn が正則化なし OT 問題の $\tau$-近似解を $O(n^2\log(n) \tau^{-3})$ 演算で計算することを意味する。丸めスキームは、2つのベクトル $\uD\in\RR^n,\vD\in\RR^m$ が与えられたとき、以下の更新を行うことからなる（\cite[Alg. 2]{altschuler2017near}）：
	$$\begin{aligned}
		\uD'& \eqdef \uD\odot\min\left(\frac{\a}{\uD\odot(\K\vD)},\ones_n\right),
		\vD' \eqdef \vD\odot \min\left(\frac{\b}{\vD\odot(\transp{\K}\uD')},\ones_n\right), \\
	\Delta_{\a}& \eqdef \a- \uD'\odot (\K \vD'), \Delta_{\b} \eqdef \b- \vD'\odot (\transp{\K} \uD), \\
	\hat\P& \eqdef \diag(\uD')\K\diag(\vD')+\Delta_{\a} \transp{(\Delta_{\b})}/\norm{\Delta_{\a}}_1.
	\end{aligned}$$
	これにより $\hat\P \in \CouplingsD(\a,\b)$ なる行列が得られ、$\hat\P$ と $\diag(\uD)\K\diag(\vD)$ の間の $1$-ノルムは $\diag(\uD)\K\diag(\vD)$ の周辺違反で制御される。すなわち
	$$ \norm{\hat\P-\diag(\uD)\K\diag(\vD)}_1 \leq \norm{\a-\uD\odot(\K\vD)}_1+\norm{\b-\vD\odot(\transp{\K}\uD)}_1.$$
この分野は依然として活発であり、\citet{pmlr-v80-dvurechensky18a}による上記の結果の最近の改良がそれを示している。
\end{rem1}
%%%%%%%%%%

%%%%%%%
\begin{rem}[Sinkhorn 反復の数値的安定性]\label{rem-stability}
注意~\ref{rem-global-conv-sinkh}および注意~\ref{rem-local-conv}で議論するように、Sinkhorn アルゴリズムの収束は $\varepsilon\rightarrow 0$ のとき劣化する。しかし数値的な実践では、この低速化がより単純な理由で実際にはほとんど観察されないことがある：カーネル $\K$ の要素の一部がメモリ上で正の数として格納するには小さすぎ、零になってしまうと、Sinkhorn アルゴリズムはしばしば終了に失敗する。これにより行列積 $\K\vD$ や $\transp{\K}\uD$ の成分がますます小さくなり零になって、式~\eqref{eq-sinkhorn}の Sinkhorn 更新において $0$ による除算が生じうる。このような問題は、乗算子 $\uD$ と $\vD$ に関する計算を対数領域で行うことで部分的に解決できる。このアプローチは注意~\ref{rem-log-sinkh}で詳細に提示され、問題~\eqref{eq-regularized-discr}の双対の直接的な解法と関連している。
\end{rem}
%%%%%%%

%%%%%%%
\begin{rem}[反復射影との関係]\label{rem-iterative-projection}
以下を定義する。
\eq{
	\Cc^1_\a \eqdef \enscond{\P}{\P\ones_m=\a}
	\qandq
	\Cc^2_\b \eqdef \enscond{\P}{\transp{\P}\ones_m=\b}
}
行および列の制約とすると、$\CouplingsD(\a,\b) = \Cc^1_\a \cap \Cc^2_\b$ である。Bregman 反復射影~\citep{bregman1967relaxation}を用いることができる。
\eql{\label{eq-kl-sinkh-proj}
	\itt{\P} \eqdef \Proj_{\Cc^1_\a}^{\KLD}(\it{\P})
	\qandq
	\ittt{\P} \eqdef \Proj_{\Cc^2_\b}^{\KLD}(\itt{\P}).
}
集合 $\Cc^1_\a$ と $\Cc^2_\b$ はアフィンであるから、これらの反復は~\eqref{eq-kl-proj}の解に収束することが知られている。\citep{bregman1967relaxation}を参照されたい。これらの反復は Sinkhorn 反復~\eqref{eq-sinkhorn}と等価である。
\eq{\label{eq-sink-matrix}\P^{(2\ell)} \eqdef \diag(\it{\uD}) \K \diag(\it{\vD}),}
と定義すると、
\begin{align*}
	\P^{(2\ell+1)} &\eqdef \diag(\itt{\uD}) \K \diag(\it{\vD}) \\
	\qandq
	\P^{(2\ell+2)} &\eqdef \diag(\itt{\uD}) \K \diag(\itt{\vD}).
\end{align*}
しかし実際には、スケーリングベクトルの操作と Gibbs カーネルとの乗算のみを必要とする~\eqref{eq-sinkhorn}を用いるべきである。これはしばしば高速化が可能である（以下の注意~\ref{rem-separable}および注意~\ref{rem-geod-heat}を参照）。
\end{rem}
%%%%%%%

%%%%%%%
\begin{rem}[近接点アルゴリズム]
正則化なし（$\epsilon=0$）の問題~\eqref{eq-mk-discr}の解を近似するために、$\KL$ 計量に対するいわゆる近接点アルゴリズムを用いて Sinkhorn アルゴリズムを反復的に使用することができる。
$F(\P) \eqdef \dotp{\P}{\pi} + \iota_{\CouplingsD(\a,\b)}(\P)$ を正則化なし目的関数とする。
$\KLD$ ダイバージェンスに対する近接点反復は $F$ の最小化子を計算し、したがって正則化なし OT 問題~\eqref{eq-mk-discr}の解を以下の反復により計算する。
\eql{\label{eq-prox-point}
	\itt{\P} \eqdef \Prox_{\frac{1}{\epsilon} F}^{\KLD}(\it{\P})
	\eqdef
	\uargmin{\P \in \RR_+^{n \times m}} \KLD(\P|\it{\P}) + \frac{1}{\epsilon} F(\P)
}
任意の $\init{\P}$ から開始する（\eqref{eq-prox-kl}も参照）。
近接点アルゴリズムは最も基本的な近接分裂法である。
ユークリッド計量に対して最初に導入され（例えば \citealt{rockafellar1976monotone}を参照）、任意の Bregman ダイバージェンスに拡張される~\citep{censor1992proximal}。したがって特にここでは $\KLD$ ダイバージェンスに対して適用できる（注意~\ref{rem-bregman}を参照）。
近接作用素は一般に閉じた形では利用できないため、何らかの形の部分反復が必要となる。
\eqref{eq-prox-point}に現れる最適化は、エントロピー正則化問題~\eqref{eq-regularized-discr}と非常に類似しており、負エントロピー $-\HD$ の代わりに相対エントロピー $\KLD(\cdot|\it{\P})$ が用いられている。
命題~\ref{prop-regularized-primal}および Sinkhorn 反復~\eqref{eq-sinkhorn}は、Gibbs カーネルを $\K=e^{-\frac{\C}{\epsilon}} \odot \it{\P} = ( e^{-\frac{\C_{i,j}}{\epsilon}} \it{\P}_{i,j} )_{i,j}$ と定義することで、このより一般的な設定にも引き継がれる。
反復~\eqref{eq-prox-point}は、各反復で Sinkhorn アルゴリズムを実行することで実装できる。
簡潔さのため $\init{\P}=\ones_n \ones_m^\top$ と仮定すると、これらの反復は以下の形を持つ。
\begin{align*}
	\itt{P} &=
		\diag(\it{\uD}) ( e^{-\frac{\C}{\epsilon}} \odot \it{\P} ) \diag(\it{\vD})\\
		&=
		\diag(\it{\uD} \odot \cdots \odot \init{\uD})
			e^{-\frac{(\ell+1)\C}{\epsilon}} \odot \it{\P} )
		\diag(\it{\vD} \odot \cdots \odot \init{\vD}).
\end{align*}
したがって近接点反復は、カーネル $e^{-\frac{\C}{\epsilon/\ell}}$ で Sinkhorn アルゴリズムを反復的に適用する。すなわち、減衰する正則化パラメータ $\epsilon/\ell$ を用いる。
この手法は、正則化に対する減衰スケジュールと Sinkhorn を組み合わせた一連の研究と密接に関連している。例えば~\citep{kosowsky1994invisible}を参照されたい。これらは、適応的なスパースグリッド上でカップリングを近似するマルチグリッド戦略と組み合わせた場合、小さな空間次元で効率的である~\citep{schmitzer2016stabilized}。
\end{rem}
%%%%%%%


%%%%%%%%%
\begin{rem}[他の正則化]
\eqref{eq-regularized-discr}のエントロピー項 $-\HD(\P)$ を任意の狭義凸ペナルティ $R(\P)$ で置き換えることが可能である。詳細は~\citep{dessein2016regularized}を参照されたい。典型的な例は $\ell^2$ ノルムの2乗であり、
\eql{\label{eq-quad-regul}
	R(\P) = \sum_{i,j} \P_{i,j}^2 + \iota_{\RR_+}(\P_{i,j});
}
\citep{essid2017quadratically}を参照されたい。
別の例は Tsallis エントロピーの族~\citep{muzellec2017tsallis}である。
ただし、ペナルティ関数が $\P$ の成分が非正の場合にも定義される場合（例えば2次正則化~\eqref{eq-quad-regul}の場合がそうである）、周辺制約 $\P\ones_m=\a$ と $\P^\top \ones_n=\b$ に加えて非負制約 $\P \geq 0$ を付け加えなければならない。実際、エントロピーを用いる場合は $\P$ の成分を正の象限にとどめる対数項を含むため非負制約を無視できるが、そうでない場合は制約の集合がアフィンでなくなり、Bregman 反復射影がもはや解に収束しなくなる。
回避策は \citeauthor{Dykstra83} のアルゴリズム~(\citeyear{Dykstra83,Dykstra85})（\citealt{bauschke-lewis}も参照）を代わりに用いることであり、\citep{2015-benamou-cisc}で詳述されている。このアルゴリズムは $R$ に付随する Bregman ダイバージェンスに従った射影を用いる。Bregman ダイバージェンスの詳細については注意~\ref{rem-bregman}を参照されたい。一般にこれらの射影は陽に計算できないという問題がある。
2乗ノルム~\eqref{eq-quad-regul}の場合、これは $(\Cc^1_\a,\Cc^2_\b)$ 上のユークリッド射影（追加の正値制約付き）の計算に対応し、単体上の射影アルゴリズム~\citep{condat2015fast}を用いて効率的に解くことができる。エントロピーに対する2次正則化の主な利点は、最適カップリングのスパースな近似を生成することであるが、これはより遅いアルゴリズムという代償を伴い、複数の optimal transport を同時に計算するために Sinkhorn ほど効率的に並列化できない（第~\ref{rem-parallel}節で議論する）。図~\ref{fig-quad-regul}はエントロピーと2次正則化による近似を対比している。
\end{rem}
%%%%%%%%%


\newcommand{\myfigRegQuad}[1]{\imgBox{\includegraphics[width=.185\linewidth,trim=65 55 45 40,clip]{quadratic-regul/quadratic-levelsets-#1}}}
\newcommand{\myfigRegEntrop}[1]{\imgBox{\includegraphics[width=.185\linewidth,trim=65 55 45 40,clip]{quadratic-regul/entropy-levelsets-#1}}}

\begin{figure}[h!]
\begin{tabular}{@{}c@{}c@{}c@{}c@{}c@{}}
\myfigRegEntrop{1} &
\myfigRegEntrop{3} &
\myfigRegEntrop{4} &
\myfigRegEntrop{7} &
\myfigRegEntrop{9} \\
{\color[rgb]{.6,.6,0} $\varepsilon=10$} &
{\color[rgb]{.45,.6,.15} $\varepsilon=1$} &
{\color[rgb]{.3,.6,.3} $\varepsilon=0.5 \cdot 10^{-1}$} &
{\color[rgb]{.15,.6,.45} $\varepsilon=10^{-1}$} &
{\color[rgb]{0,.6,.6} $\varepsilon=10^{-3}$} \\
\myfigRegQuad{1} &
\myfigRegQuad{3} &
\myfigRegQuad{4} &
\myfigRegQuad{7} &
\myfigRegQuad{9} \\
{\color[rgb]{.6,.6,0} $\varepsilon=5 \cdot 10^3$} &
{\color[rgb]{.45,.6,.15}$\varepsilon=10^3$} &
{\color[rgb]{.3,.6,.3} $\varepsilon=10^2$} &
{\color[rgb]{.15,.6,.45} $\varepsilon=10$} &
{\color[rgb]{0,.6,.6} $\varepsilon=1$}
\end{tabular}
\caption{\label{fig-quad-regul}
エントロピー正則化 $R=-\HD$（上段）と2次正則化 $R=\norm{\cdot}^2+\iota_{\RR_+}$（下段）の比較。周辺分布 $(\al,\be)$ は図~\ref{fig-entropic-copula}と同じである。
}
\end{figure}



%%%%%%%%%%
\begin{rem1}{重心射影}\label{rem-barycenric-proj}
	注意~\ref{rem:entrop-reg}の設定、すなわちエントロピー正則化を用いて離散測度間の OT を近似する場合を再び考える。\eqref{eq-mk-discr}の Kantorovich 定式化とそのエントロピー正則化~\eqref{eq-regularized-discr}の両方がカップリング $\P \in \CouplingsD(\a,\b)$ を与える。$\Y=\RR^\dim$ の場合に輸送写像 $\T : \X \rightarrow \Y$ を定義するために、いわゆる重心射影写像を定義できる。
	\eql{\label{eq-baryproj}
		\T : x_i \in \X \longmapsto \frac{1}{\a_i} \sum_j \P_{i,j} y_j \in \Y,
	}
	ここで入力測度は~\eqref{eq-pair-discr}の形の離散測度である。この写像は $\al$ の台の点 $(x_i)_i$ に対してのみ定義されることに注意されたい。
	$\T$ が置換行列である場合（命題~\ref{prop-matching-kanto}で詳述されるように）、$\T$ は Monge 写像に等しく、$\varepsilon \rightarrow 0$ のとき重心射影はその写像が一意であれば漸進的にその写像に収束する。
	任意の（必ずしも離散でない）測度に対しては、\eqref{eq-mk-generic}またはその正則化版~\eqref{eq-entropic-generic}を解くことでカップリング $\pi \in \Couplings(\al,\be)$ が定義される。このカップリング $\pi$ は常に $\al \otimes \be$ に関する密度 $\frac{\d\pi(x,y)}{\d\al(x)\d\be(y)}$ を持つことに注意されたい。したがって写像は以下の公式で取り出すことができる。
	\eql{\label{eq-bary-proj}
		\T : x \in \X \longmapsto \int_\Yy y \frac{\d\pi(x,y)}{\d\al(x)\d\be(y)} \d\be(y).
	}
	$\varepsilon=0$ の場合に $\pi$ が Monge 写像のグラフ上に台を持つとき（注意~\ref{rem-exist-mongemap}を参照）、$\varepsilon>0$ を用いるとこの写像の滑らかな近似が得られる。
	このような重心射影は、OT の Monge 写像をイメージングの問題に適用する際に有用である。色変換への応用については図~\ref{fig-colors}を参照されたい。また、Wasserstein 計量を備えた確率測度の空間における主測地線の近似の計算にも用いられている。\citep{SeguyCuturi}を参照されたい。
\end{rem1}
%%%%%%%%%



%%%%%%%%%%
\begin{rem1}{Hilbert 計量}
\citep{franklin1989scaling}が最初に説明したように、Sinkhorn の大域的収束解析は $\RR_{+,*}^n$（正ベクトル）上の Hilbert 射影計量を用いると大幅に簡略化される。これは以下で定義される。
\eq{
	\foralls (\uD,\uD') \in (\RR_{+,*}^n)^2, \quad
	\Hilbert(\uD,\uD') \eqdef \log \umax{i,j} \frac{ \uD_i \uD_{j}' }{ \uD_{j} \uD_{i}'  }.
}
これは射影錐 $\RR_{+,*}^n/\sim$ 上の距離であることが示される。ここで $\uD \sim \uD'$ は $\exists r>0, \uD=r\uD'$ を意味する（ベクトルがスケーリングを除いて等しいことから「射影的」という名称がつく）。
すなわち、$\Hilbert$ は三角不等式を満たし、$\Hilbert(\uD,\uD')=0$ であることと $\uD \sim \uD'$ であることは同値である。
これは有界開凸集合上の Hilbert のもとの距離~\citep{hilbert1895gerade}の射影版である。
射影錐 $\RR_{+,*}^n/\sim$ はこの距離に関して完備距離空間である。
対数変数変換により、正の錐の射線上の Hilbert 計量は変動セミノルムに等距離的である（これは加法定数を除いて定義されるベクトル間のノルムである）。
\eql{\label{eq-hilbert-var}
	\Hilbert(\uD,\uD') = \norm{\log(\uD)-\log(\uD')}_{\text{var}}
}
\eq{
	\qwhereq
	\norm{\fD}_{\text{var}} \eqdef (\max_i \fD_i) - (\min_i \fD_i).
}
この変動セミノルムは $\ell^\infty$ ノルムと密接に関連しており、常に $\norm{\fD}_{\text{var}}  \leq 2 \norm{\fD}_\infty$ が成り立つ。ある固定された $i$ に対して $\fD_i=0$ を課すと、逆の不等式 $\norm{\fD}_\infty \leq \norm{\fD}_{\text{var}}$ も成り立つ。これらの評価は Sinkhorn の収束解析（以下の注意~\ref{rem-global-conv-sinkh}を参照）において特に有用である。なぜなら、~\eqref{eq-dualsinkhorn-constraints2}を解く双対変数 $\fD=\log(\uD)$ は加法定数を除いて定義されるため、ある $i$ に対して $\fD_i=0$ を課すことができるからである。
Hilbert 計量は \citep{birkhoff1957extensions} と \citep{samelson1957perron} によって独立に導入された。彼らは以下の基本定理を証明した。これは正行列が正ベクトルの錐上で厳密な縮小であることを示している。


\begin{thm}\label{thm-birkoff}
	$\K \in \RR_{+,*}^{n \times m}$ とする。すると $(\vD,\vD') \in (\RR_{+,*}^m)^2$ に対して
	\eq{
		\Hilbert(\K \vD,\K \vD') \leq \la(\K) \Hilbert(\vD,\vD'),
		\text{ where }
		\choice{
			\la(\K) \eqdef \frac{ \sqrt{\eta(\K)}-1 }{ \sqrt{\eta(\K)}+1 } < 1, \\
			\eta(\K) \eqdef \umax{i,j,k,\ell} \frac{ \K_{i,k} \K_{j,\ell} }{ \K_{j,k} \K_{i,\ell} }.
		}
	}
\end{thm}

図~\ref{fig-h-m}はこの定理を例示している。
\end{rem1}


\begin{figure}[h!]
\centering
\includegraphics[width=.22\linewidth]{hilbert-metric/hilbert-metric} \qquad
\includegraphics[width=.7\linewidth]{hilbert-metric/cones}\\
\caption{\label{fig-h-m}
左：Hilbert 計量 $\Hilbert$ は錐内の射線上の距離である（ここでは正ベクトル）。
右：正行列 $\K$ の反復により誘導される縮小の可視化。
}
\end{figure}


%%%%%%%%%%
\begin{rem1}{Perron--Frobenius}
定理~\ref{thm-birkoff}の典型的な応用は、Perron--Frobenius の定理の定量的証明を与えることである。注意~\ref{rem-local-conv}で説明されるように、これは Sinkhorn 反復の局所的線形化と関連している。$\K^\top \ones_n = \ones_n$ を満たす行列 $\K \in \RR_+^{n \times n}$ は $\simplex_n$ を $\simplex_n$ に写す。さらに $\K>0$ であれば、定理~\ref{thm-birkoff}により計量 $\Hilbert$ に関して厳密に縮小的であり、$\K p^\star=p^\star$ を満たす一意の不変確率分布 $p^\star \in \simplex_n$ が存在する。さらに、任意の $p_0 \in \simplex_n$ に対して $\Hilbert(\K^\ell p_0,p^\star) \leq \la(\K)^\ell \Hilbert(p_0,p^\star)$ が成り立つ。\ie 行列の反復の $p^\star$ への線形収束が得られる。これは図~\ref{fig-perron}に示されている。
\end{rem1}
%%%%%%%%%%

\begin{figure}[h!]
\centering
\begin{tabular}{@{}c@{}c@{}c@{\hspace{10mm}}c@{}}
\includegraphics[width=.22\linewidth]{hilbert-metric/simplex-0}&
\includegraphics[width=.22\linewidth]{hilbert-metric/simplex-1}&
\includegraphics[width=.22\linewidth]{hilbert-metric/simplex-2}&
\includegraphics[width=.22\linewidth]{hilbert-metric/simplex-evol}\\
$\simplex_3$ & $\K \simplex_3$ & $\K^2 \simplex_3$ &  $\{ \K^\ell \simplex_3 \}_\ell$
\end{tabular}
\caption{\label{fig-perron}
$\K^\top \ones_3=\ones_3$ を満たす $\K \in \RR_{+,*}^{3 \times 3}$ の不変確率分布への $\K^\ell \simplex_3 \rightarrow \{p^\star\}$ の変化。
}
\end{figure}




%%%%%%%%%%
\begin{rem11}{大域的収束}{rem-global-conv-sinkh}
以下の定理は \citep{franklin1989scaling} により証明され、定理~\ref{thm-birkoff}を用いて Sinkhorn 反復の線形収束を示す。

\begin{thm}
	$(\it{\uD},\it{\vD}) \rightarrow (\uD^\star,\vD^\star)$ が成り立ち、
	\eql{\label{eq-convlin-sinkh}
		\Hilbert(\it{\uD}, \uD^\star) = O(\la(\K)^{2\ell}), \quad
		\Hilbert(\it{\vD}, \vD^\star) = O(\la(\K)^{2\ell}).
	}
	さらに以下も成り立つ。
	\eql{\label{eq-convsinkh-control}
		\begin{split}
		\Hilbert(\it{\uD}, \uD^\star) &\leq \frac{\Hilbert( \it{\P}\ones_m,\a )}{1-\la(\K)^2}, \\
		\Hilbert(\it{\vD}, \vD^\star) &\leq \frac{\Hilbert( \P^{(\ell),\top} \ones_n,\b )}{1-\la(\K)^2},
		\end{split}
	}
	ここで $\it{\P} \eqdef \diag(\it{\uD}) \K \diag(\it{\vD})$ とした。最後に以下が成り立つ。
	\eql{\label{eq-convlin-sinkh-prim}
		\|\log(\it{\P}) - \log(\P^\star)\|_\infty \leq \Hilbert(\it{\uD}, \uD^\star) + \Hilbert(\it{\vD}, \vD^\star),
	}
	ここで $\P^\star$ は~\eqref{eq-regularized-discr}の一意の解である。
\end{thm}

\begin{proof}
	任意の $(\vD,\vD') \in (\RR_{+,*}^m)^2$ に対して以下が成り立つことに注意する。
	\eq{
		\Hilbert(\vD,\vD') = \Hilbert(\vD/\vD',\ones_m) = \Hilbert(\ones_m/\vD,\ones_m/\vD').
	}
	これにより
	\begin{align*}
		\Hilbert(\itt{\uD},\uD^\star) &= \Hilbert\pa{ \frac{\a}{\K \it{\vD}}, \frac{\a}{\K \vD^\star} } \\
		&= \Hilbert( \K \it{\vD}, \K \vD^\star ) \leq \la(\K) \Hilbert( \it{\vD}, \vD^\star ),
	\end{align*}
	ここで定理~\ref{thm-birkoff}を用いた。これにより~\eqref{eq-convlin-sinkh}が示される。三角不等式を用いると、
	\begin{align*}
		\Hilbert(\it{\uD},\uD^\star) &\leq \Hilbert(\itt{\uD},\it{\uD}) + \Hilbert(\itt{\uD},\uD^\star) \\
		&\leq \Hilbert\pa{ \frac{\a}{\K \it{\vD}},\it{\uD} } + \la(\K)^2 \Hilbert(\it{\uD},\uD^\star) \\
		&= \Hilbert\pa{ \a,\it{\uD} \odot  ( \K \it{\vD} ) } + \la(\K)^2 \Hilbert(\it{\uD},\uD^\star),
	\end{align*}
	が得られ、$\it{\uD} \odot  ( \K \it{\vD} ) = \it{\P}\ones_m$ であるから~\eqref{eq-convsinkh-control}の最初の部分が得られる（2番目も同様）。
	\eqref{eq-convlin-sinkh-prim}の証明は~\citep[Lem. 3]{franklin1989scaling}に従う。
\end{proof}

評価~\eqref{eq-convsinkh-control}は、周辺制約違反に関する誤差指標、例えば $\| \it{\P} \ones_m - \a \|_1$ や $\|\transp{\it{\P}} \ones_n - \b \|_1$ が収束を監視するための有用な停止基準であることを示している。
\eqref{eq-hilbert-var}のおかげで、スケーリング変数 $(\it{\uD},\it{\vD})$ に関するこれらの Hilbert 計量の収束率は、双対変数 $(\it{\fD},\it{\gD}) \eqdef (\varepsilon\log(\it{\uD}),\varepsilon\log(\it{\vD}))$ の変動ノルム $\norm{\cdot}_{\text{var}}$ に関する線形収束率を与える。

図~\ref{fig-sinkhorn-convergence}の下段は、制約違反に関するこの線形収束率を示し、$\varepsilon\rightarrow 0$ のときこの収束率がどのように劣化するかを示している。
これらの結果は~\citep{franklin1989scaling}で証明され、非線形 Perron--Frobenius 理論~\citep{lemmens2012nonlinear}と密接に関連している。Perron--Frobenius 理論は反復の線形化に対応する。\eqref{eq-linearized-sinkh}を参照されたい。この収束解析は \citep{linial1998deterministic} により拡張され、Sinkhorn の各反復がスケーリングされたカップリング行列の permanent を増加させることが示されている。
\end{rem11}
%%%%%%%%%%



%%%%%%%%%%
\begin{rem1}{局所的収束}\label{rem-local-conv}
大域的線形収束率~\eqref{eq-convlin-sinkh-prim}は、$\X=\Y=\RR^d$ で $\varepsilon=0$ のとき Monge 写像が存在する場合（注意~\ref{rem-monge2}を参照）のように、しばしばかなり悲観的である。大域的収束率は、コスト行列 $\C$ がランダムに近いより困難な状況ではむしろ鋭く、これらの場合、収束率は $\varepsilon$ に関して指数的に悪化する。$1-\la(\K) \sim e^{-1/\varepsilon}$ である。
より精密な漸近解析（\eg 高精度の解に関心があり多数の反復を行う場合）を得るには、通常、局所的収束率を研究する。
Sinkhorn の更新を固定点写像 $\itt{\fD} = \Phi( \it{\fD} )$ の反復として書くことができ、ここで
\eq{
	\Phi \eqdef \Phi_2 \odot \Phi_1
	\qwhereq
	\choice{
		\Phi_1(\fD) = \varepsilon \log \transp{\K} ( e^{\fD/\varepsilon} ) - \log(\b), \\
		\Phi_2(\gD) = \varepsilon \log \K ( e^{\gD/\varepsilon} ) - \log(\a).
	}
}
~\eqref{eq-dual-formulation}を解く最適な $(\fD,\gD)$ に対して、$\P = \diag(e^{\fD/\varepsilon}) \K \diag(e^{\gD/\varepsilon})$ を~\eqref{eq-regularized-discr}を解く最適カップリングとすると、以下の Jacobi 行列が得られる：
\eql{\label{eq-linearized-sinkh}
	\partial \Phi(\fD) = \diag(\a)^{-1} \odot \P \odot \diag(\b)^{-1} \odot \transp{\P}.
}
この Jacobi 行列は $\partial \Phi(\fD) \ones_n = \ones_n$ を満たす正行列であり、したがって Perron--Frobenius の定理により、固有値 $1$ に対する唯一の支配固有ベクトル $\ones_m$ を持つ。$\fD$ は定数を除いて定義されるので、実際には局所的な線形収束率を支配するのは第2固有値 $1-\kappa<1$ であり、$\ell$ が十分大きいとき
\eq{
	\|\it{\fD}-\fD \| = O( (1-\kappa)^\ell ).
}
数値的に、「単純な場合」（$\varepsilon=0$ のとき滑らかな Monge 写像が存在する場合など）では、この収束率は $\kappa \sim \varepsilon$ のようにスケールする。
詳細については、双確率的（割当）の場合の~\citep{knight2008sinkhorn}を参照されたい。
\end{rem1}
%%%%%%%%%%






%%%%%%%%%%%%%%%%%%%%%%%%%%%%%%%%%%%%%%%%%%%%%%%%%%%%%%%%%%%%%%%%%%%%%%%%%%%%%
\section{Sinkhorn 反復の高速化}

Sinkhorn 反復の主な計算のボトルネックは、カーネル $\K$ および $\K^\top$ に対するベクトル-行列積であり、素朴に実装すると計算量は $O(nm)$ である。以下では、計算量を大幅に改善できるいくつかの重要な場合を詳述する。

%%%%%%%
\begin{rem}[並列・GPU 向け計算]\label{rem-parallel} Sinkhorn アルゴリズムの単純さは、ヒストグラムのペア間の複数の正則化 Wasserstein 距離を同時に計算する非常に効率的なアプローチを与える。
$N$ を整数とし、$\a_1,\ldots,\a_N$ を $\simplex_n$ のヒストグラム、$\b_1,\ldots,\b_N$ を $\Sigma_m$ のヒストグラムとする。
$N$ 個の近似距離 $\MKD_\C^\varepsilon(\a_1,\b_1),\dots,\MKD_\C^\varepsilon(\a_N,\b_N)$ をすべて計算することを目指す。
この場合、$\A=[\a_1,\dots,\a_N]$ および $\B=[\b_1,\dots,\b_N]$ をすべてのヒストグラムを格納する $n\times N$ および $m\times N$ 行列とすると、これら $N$ ペアすべての Sinkhorn 反復は以下のように並列に実行できる。
\eql{\label{eq-sinkhorn-par}
	\itt{\UD} \eqdef \frac{\A}{\K \it{\VD}}
	\qandq
	\itt{\VD} \eqdef \frac{\B}{\transp{\K}\itt{\UD}},
}
初期値は $\init{\VD} = \ones_{m \times N}$ である。ここで $\frac{\cdot}{\cdot}$ は行列の要素ごとの除算に対応する。
$\VD$ と $\UD$ の収束時に、正則化距離の（行）ベクトルが以下に簡約されることが確認できる。
$$ \transp{\ones_n} (    \UD\odot \log \UD \odot (   (\K\odot \C) \VD) +  \UD\odot(   (\K\odot \C)(\VD\odot \log \VD)    ) ) \in\RR^{N}.$$
式~\eqref{eq-sinkhorn}で記述される基本的な Sinkhorn 反復は本質的に GPU 向きである。行列-ベクトル積のみからなるためであり、これは例えば \citet{slomp2011gpu} でマッチング問題を解くために利用された。しかし、式~\eqref{eq-sinkhorn-par}で示される行列-行列演算は並列性のさらに良い機会を提供し、大規模なヒストグラム間の OT 距離を計算する際の Sinkhorn アルゴリズムの成功を説明している。
\end{rem}
%%%%%%%

%%%%%%%%%
\begin{rem}[分離可能カーネルによる高速化]\label{rem-separable} 本節では、各 Sinkhorn 反復の計算量を大幅に削減できる重要な特殊ケースを考える。この特殊ケースは、コスト行列で考慮される各インデックス $i$ と $j$ が $d$ 個の有限集合 $\range{n_1},\dots,\range{n_d}$ の直積から取った $d$-組として記述できる場合に生じる。
$$i=(i_k)_{k=1}^d, j=(j_k)_{k=1}^d \in \range{n_1}\times\dots\times\range{n_d}.$$
この設定で、インデックス $i$ と $j$ の間のコスト $\C_{ij}$ がこれらのサブインデックスに沿って加法的である場合、すなわちそれぞれサイズ $n_1\times n_,\dots,n_d\times n_d$ の $d$ 個の行列 $\C^1,\dots,\C^d$ が存在して
$$\C_{ij} = \sum_{k=1}^d \C^k_{i_k,j_k},$$
が成り立つ場合、直接的な帰結として、Sinkhorn 反復に現れるカーネルは分離可能な乗法構造を持つ。
\eql{\label{eq-speedup-separable}
\K_{i,j} = \prod_{k=1}^d \K^k_{i_k,j_k}.
}
このような分離可能な乗法構造は $\K \uD$ の非常に高速な（厳密な）評価を可能にする。実際、$n=\prod_k n_k$ は通常次元 $d$ に関して指数的であるため、$\K$ をサイズ $n\times n$ の行列としてインスタンス化する代わりに（これは法外なサイズとなる）、$\uD$ の各「スライス」に沿って $\K^k$ を適用するだけで $\K\uD$ を復元できる。$n=m$ の場合、計算量は $O(n^2)$ の代わりに $O(n^{1+1/d})$ に削減される。

この高速化の重要な例は、$\Xx=\Yy=[0,1]^d$ で、ground cost が $q$-ノルムの $q$ 乗、$$\c(x,y)=\norm{x-y}^q_q=\sum_{i=1}^d |x_i-y_i|^q, \; q>0;$$ であり、空間が正則格子を用いて離散化され、点 $x_i = (i_1/n_1,\ldots,i_d/n_d)$（$i=(i_1,\dots,i_d)\in\range{n_1}\times \dots\times\range{n_d}$）のみが考慮される場合に生じる。この場合、$\K$ による乗算は、各1次元 $n_k\times n_k$ 畳み込み行列 $$\K^k = \begin{bmatrix}\exp(-\left|\frac{r-s}{n_k}\right|^q/ \varepsilon)\end{bmatrix}_{1 \leq r,s\leq n_k}$$ を、第一の次元が $k$ 番目のインデックスの集合と一致するように置換されたテンソルとして再形成された $\uD$ に適用することで、より効率的に実行できる。
例えば、$d=2$（平面の場合）かつ $q=2$（2-Wasserstein、Gauss 畳み込みに帰着）の場合、ヒストグラム $\a$ したがって Sinkhorn の乗算子 $\uD$ は $n_1\times n_2$ 行列としてインスタンス化できる。乗算子 $\uD$ が長さ $n_1n_2$ のベクトルではなく $n_1\times n_2$ 行列として再形成されていることを強調するために $\mathbf{U}$ と書く。すると、素朴な実装では $(n_1 n_2)^2$ 演算を必要とする $\K\uD$ の計算は、2つの1次元畳み込みを個別に適用することで、$$(\K^2(\K^1 \mathbf{U})^T)^T=\K^1\mathbf{U}\K^2,$$ として $n_1^2n_2^2$ 演算の代わりに $(n_1^2)n_2+n_1(n_2^2)$ 演算で $n_1\times n_2$ 行列を得ることができる。この例は上で与えた指数 $(1+1/d)$ と一致することに注意されたい。より大きな $d$ では、$\uD$（適切なサイズのテンソルとして再形成）の各スライスにこれらと全く同じ1次元畳み込みを適用する必要があり、これは GPU 上で非常に効率的な操作である。

この重要な観察は、\citep{2015-solomon-siggraph,2016-bonneel-barycoord}で強調されているように、2次元および3次元の形状データに optimal transport を適用する際の実用的成功の多くを支えている。そこでは $200^3=8 \times 10^6$ もの大きさの格子上に台を持つ分布が扱われている。
\end{rem}
%%%%%%%%%

%%%%%%%%%
\begin{rem}[近似畳み込み]\label{rem-convol-sinkh}
Sinkhorn 反復~\eqref{eq-sinkhorn}の主な計算のボトルネックは、$\K$ またはその随伴によるベクトルの乗算にある。分離可能性~\eqref{eq-speedup-separable}の利用に加えて、カーネルの他の特殊構造を活用することも可能である。最も単純なケースは、平行移動不変カーネル $\K_{i,j} = k_{i-j}$ の場合であり、ユークリッド空間 $\X=\RR^\dim$ における固定された一様格子上で測度を離散化する場合が典型的である。このとき $\K \vD = k \star \vD$ は畳み込みであり、畳み込みをほぼ線形時間で近似するいくつかのアルゴリズムが存在する。
最も一般的なのは Fourier 変換 $\Ff$ によるものであり、簡潔さのため周期的境界条件を仮定すると $\Ff(k \star \vD)=\Ff(k) \odot \Ff(\vD)$ が成り立つ。しかし、これは不安定な計算につながり、小さな $\varepsilon$ に対してはしばしば許容されない。
別の一般的な高速化法は、Deriche フィルタリング法 \citet{deriche1993recursively} などを用いた自己回帰フィルタの連続適用により畳み込みを近似することである。様々な高速フィルタリング法の比較については~\citep{getreuer2013survey}を参照されたい。
\end{rem}
%%%%%%%%%

%%%%%%%%%
\begin{rem}[熱における測地線近似]\label{rem-geod-heat}
非平面領域では、カーネル $\K$ は畳み込みではないが、コストが $\C_{i,j}=d_\Mm(x_i,y_j)^p$（ここで $d_\Mm$ は曲面 $\Mm$（またはより一般的な多様体）上の測地線距離）の場合、$\K=e^{-\frac{d_\Mm}{\varepsilon}}$ のベクトルへの適用を高速に近似することも可能である。実際、\citeauthor{varadhan-1967}の公式~\citeyearpar{varadhan-1967}は、このカーネルが Laplacian カーネル（$p=1$ の場合）および熱核（$p=2$ の場合）に近いことを主張する。
Varadhan の第1公式は以下を述べる。
\eql{\label{eq-varadhan-1}
	-\frac{\sqrt{t}}{2} \log( \Pp_t(x,y) ) = d_\Mm(x,y) + o(t)
	\qwhereq
	\Pp_t \eqdef (\Id - t\De_{\Mm})^{-1},
}
ここで $\De_{\Mm}$ は多様体 $\Mm$ に付随する Laplace--Beltrami 作用素（半負定値）であり、$\Pp_t$ は積分核であり、$g = \int_\Mm \Pp_t(x,y) f(y) \d y$ は $g - t \De_{\Mm} g = f$ の解である。
Varadhan の第2公式は以下を述べる。
\eql{\label{eq-varadhan-2}
	\sqrt{ -4 t \log( \Hh_t(x,y) ) } = d_\Mm(x,y) + o(t),
}
ここで $\Hh_t$ は、$g_t = \int_\Mm \Hh_t(x,y) f(y) \d y$ が時刻 $t$ における熱方程式
\eq{
	\frac{\partial g_t(x)}{\partial t} = (\Delta_\Mm g_t)(x).
}
の解となるように定義される積分核である。
公式~\eqref{eq-varadhan-1}および~\eqref{eq-varadhan-2}における収束はコンパクト多様体上で一様である。
数値的には、領域 $\Mm$ は離散化され（例えば有限要素を用いて）、$\De_\Mm$ は離散 Laplacian 行列 $L$ で近似される。典型的な例は区分線形有限要素を用いる場合であり、$L$ は有名な余接 Laplacian である（この構成の詳細については~\citep{botsch-2010}を参照）。
これらの公式は Gibbs カーネル $\K_{i,j}=e^{-\frac{d(x_i,y_j)^p}{\varepsilon}}$ による乗算を効率的に近似するために用いることができる。
式~\eqref{eq-varadhan-1}は、$p=1$ の場合に $\varepsilon=\frac{\sqrt{t}}{2}$ とし、$\K$ による乗算を $(\Id-t L)^{-1}$ による乗算で置き換えることを示唆する。これは正定値対称線形系の解法を必要とする。
式~\eqref{eq-varadhan-2}は、熱流の安定的な解法のための $R$ ステップの陰的 Euler 法と組み合わせて、$p=2$ の場合に $4 t=\varepsilon$ として $\K$ による乗算を $(\Id-\frac{t}{R} L)^{-R}$ による乗算で置き換えることを示唆する。これは $R$ 回の線形系の解法を必要とする。幸い、これらの線形系は各 Sinkhorn 反復で解かれるべきものであるから、スパースな Cholesky 分解を事前に計算することで効率的に解くことができる。行列の行と列の並べ替え~\citep{george1989evolution}を行うことで、2次元多様体に対してほぼ線形のスパース性が得られ、各 Sinkhorn 反復は線形の計算量を持つ（性能は多様体の次元とともに劣化する）。
測地線距離の近似に Varadhan の公式を用いることは~\citep{Crane2013}で最初に提案され、Sinkhorn 反復と組み合わせた使用は~\citep{2015-solomon-siggraph}で提案された。
\end{rem}
%%%%%%%%%


%%%%%%%%%%
\begin{rem}[外挿加速]
	Sinkhorn アルゴリズムは固定点アルゴリズムであるから（注意~\ref{rem-local-conv}で示されたように）、標準的な線形または非線形外挿スキームを用いて、解の近傍における固定点写像の条件数を改善し、線形収束率を向上させることができる。
	これは逐次過緩和法（例えば~\citep{hadjidimos2000successive}を参照）と類似しており、局所的な線形収束率は $O((1-\kappa)^\ell)$ から $O((1-\sqrt{\kappa})^\ell)$ に改善される（ここで $\kappa>0$。注意~\ref{rem-local-conv}を参照）。
	詳細については~\citep{2016-peyre-qot}を参照されたい。
\end{rem}
%%%%%%%%%%



%%%%%%%%%%%%%%%%%%%%%%%%%%%%%%%%%%%%%%%%%%%%
\section{安定性と対数領域での計算}

注意~\ref{rem-stability}で簡単に述べたように、Sinkhorn アルゴリズムは正則化パラメータ $\epsilon$ がコスト行列 $\C$ の成分と比較して小さい場合に数値的なオーバーフローに悩まされる。この懸念は、対数領域で計算を行うことである程度緩和できる。このアプローチの妥当性は、~\eqref{eq-regularized-discr}に付随する双対問題を考えることでより明確になる。対数領域での計算はそこに自然に現れるものである。

\begin{prop}
以下が成り立つ。
\eql{\label{eq-dual-formulation}
	\MKD_\C^\varepsilon(\a,\b) = \umax{\fD \in \RR^n,\gD \in \RR^m}
		 \dotp{\fD}{\a} + \dotp{\gD}{\b}
		- \varepsilon \dotp{e^{\fD/\varepsilon} }{ \K e^{\gD/\varepsilon}}.
}
最適な $(\fD,\gD)$ は~\eqref{eq-scaling-form}に現れるスケーリング $(\uD,\vD)$ と以下の関係を通じて結びつく。
\eql{\label{eq-entropy-pd}
	(\uD,\vD)=(e^{\fD/\varepsilon},e^{\gD/\varepsilon}).
}
\end{prop}

\begin{proof}
命題~\ref{prop-regularized-primal}の証明の最後の部分から出発する。そこでは最適な主解 $\P$ と周辺制約に対する双対乗数 $\fD$ および $\gD$ が以下のように結びつけられている。
\eq{
	\label{eq-prim-dual-ent}\P_{i,j}=e^{\fD_i/\varepsilon}e^{-\C_{i,j}/\varepsilon}e^{\gD_j/\varepsilon}.
}
式~\eqref{eq-sinkhorn-lagrangian}の Lagrangian $\Lag(\P,\fD,\gD)$ に、最適な $\P$ を $\fD$ と $\gD$ の関数として代入すると、Lagrange 双対関数は以下に等しくなる。
\eql{\label{eq-lagrange-dual-naive}
	\fD,\gD \mapsto \dotp{e^{\fD/\varepsilon}}{\left(\K\odot \C\right)e^{\gD/\varepsilon}} - \varepsilon \HD(\diag(e^{\fD/\varepsilon}) \K \diag(e^{\gD/\varepsilon})).
}
$\varepsilon$ でスケールされた $\P$ の負エントロピー、すなわち $\varepsilon \dotp{\P}{\log \P - \ones_{n\times m}}$ は $\fD,\gD, \C$ の関数として陽に表すことができる。
\begin{align*}
&\dotp{\diag(e^{\fD/\varepsilon}) \K \diag(e^{\gD/\varepsilon})}{\fD\transp{\ones_m}+\ones_n\transp{\gD}-\C-\varepsilon\ones_{n\times m}}\\
&= -\dotp{e^{\fD/\varepsilon}}{\left(\K\odot \C\right)e^{\gD/\varepsilon}} + \dotp{\fD}{\a}+ \dotp{\gD}{\b}-\varepsilon \dotp{e^{\fD/\varepsilon}}{\K e^{\gD/\varepsilon}};
\end{align*}
したがって、\eqref{eq-lagrange-dual-naive}の第1項は上のエントロピーの第1項と相殺する。残りの項が~\eqref{eq-dual-formulation}に現れるものである。
\end{proof}



%%%%%%%%%%
\begin{rem}[双対問題に対するブロック座標上昇としての Sinkhorn]\label{rem-sinkh-block}
制約なし最大化問題~\eqref{eq-dual-formulation}を解く単純なアプローチは、厳密な\emph{ブロック座標上昇}戦略を用いることである。すなわち、\eqref{eq-dual-formulation}の目的関数のこれらの変数に関するそれぞれの勾配を打ち消すように $\fD$ と $\gD$ を交互に更新する。実際、\eqref{eq-dual-formulation}の目的関数を $Q(\fD,\gD)$ と書くと、いくつかの基本的な計算の後に以下が得られることに気づく。
\begin{align}
	\label{eq-dualupadate-sinkh-1}\nabla|_\fD\, Q(\fD,\gD) &=  \a - e^{\fD/\varepsilon}\odot \left(\K e^{\gD/\varepsilon}\right),\\
	\label{eq-dualupadate-sinkh-2}\nabla|_\gD\, Q(\fD,\gD) &=  \b - e^{\gD/\varepsilon}\odot \left(\transp{\K} e^{\fD/\varepsilon}\right).
\end{align}
したがってブロック座標上昇は、任意の初期値 $\init{\gD}$ から出発して、$l\geq 0$ に対し以下の更新を逐次適用することで閉じた形で実装できる：
\begin{align}
	\label{eq-slse-sinkh-1}\itt{\fD} &= \varepsilon \log \a -\varepsilon\log\left(\K e^{\it{\gD}/\varepsilon}\right), \\
	\label{eq-slse-sinkh-2}\itt{\gD} &= \varepsilon \log \b- \varepsilon\log\left(\transp{\K} e^{\itt{\fD}/\varepsilon}\right).
\end{align}
このような反復は、~\eqref{eq-entropy-pd}で強調された主双対関係を考慮すると、Sinkhorn 反復~\eqref{eq-sinkhorn}と数学的に等価である。実際、各反復において以下が復元される。
\eq{
	(\it{\fD},\it{\gD}) = \varepsilon ( \log(\it{\uD}), \log(\it{\vD}) ).
}
\end{rem}
%%%%%%%%%%

%%%%%%%%%%
\begin{rem}[ソフト最小値による書き換え]
反復~\eqref{eq-slse-sinkh-1}および~\eqref{eq-slse-sinkh-2}は、以下の記法を用いた別の解釈を与えることができる。実数のベクトル $\z$ が与えられたとき、その座標の\emph{ソフト最小値}を $\smine\z$ と書く。すなわち
\eql{\label{eq-defn-softmin}
	\smine\z= -\varepsilon \log \sum_i e^{-\z_i/\varepsilon}.
}
$\varepsilon\rightarrow 0$ のとき、任意のベクトル $\z$ に対して $\smine(\z)$ は $\min\,\z$ に収束することに注意されたい。実際、$\smine$ は図~\ref{fig-softmin}に示すように $\min$ 関数の微分可能な近似と解釈できる。

\begin{figure}[h!]
\centering
\begin{tabular}{@{}c@{\hspace{1mm}}c@{\hspace{1mm}}c@{\hspace{1mm}}c@{\hspace{1mm}}c@{}}

\includegraphics[width=.19\linewidth]{soft-min/lse-1}&
\includegraphics[width=.19\linewidth]{soft-min/lse-05}&
\includegraphics[width=.19\linewidth]{soft-min/lse-01}&
\includegraphics[width=.19\linewidth]{soft-min/lse-001}&
\includegraphics[width=.19\linewidth]{soft-min/lse-0001} \\
$\varepsilon=1$ &
$\varepsilon=0.5$ &
$\varepsilon=10^{-1}$ &
$\varepsilon=10^{-2}$ &
$\varepsilon=10^{-3}$
\end{tabular}
\caption{\label{fig-softmin}
$\varepsilon$ を変化させた場合の2次元における関数 $\smine(\z)$（$\z \in \RR^2$）の表示。
}
\end{figure}

この記法を用いると、式~\eqref{eq-slse-sinkh-1}および~\eqref{eq-slse-sinkh-2}は以下のように書き直せる。
\begin{align}
\label{eq-slse-smin-1}(\itt{\fD})_i &= \smine\, (\C_{ij}-\it{\gD}_j)_j + \varepsilon \log \a_i, \\
\label{eq-slse-smin-2}(\itt{\gD})_j &= \smine\, (\C_{ij}-\it{\fD}_i)_i + \varepsilon \log \b_j.
\end{align}
ここで項 $\smine\, (\C_{ij}-\it{\gD}_j)_j$ は、行列 $(\C-\ones_n (\it{\gD})^\top )$ の $j$ 列のすべての値のソフト最小値を表す。記法を簡略化するために、行列を入力として受け取り、その列または行のソフト最小値の列ベクトルを出力する演算子を導入する。すなわち、任意の行列 $A\in\RR^{n\times m}$ に対して以下を定義する。
\begin{align*}
	\SMINEr(\A) \eqdef \pa{\smine \pa{\A_{i,j}}_j}_i\in\RR^n,\\
	\SMINEc(\A) \eqdef \pa{\smine \pa{\A_{i,j}}_i}_j\in\RR^m.
\end{align*}
これらの演算は第~\ref{sec-semi-discr-entropy}節で導入されるエントロピー $\c$-変換と等価であることに注意されたい（特に~\eqref{eq-sinkh-c-transf}を参照）。
この記法を用いると、Sinkhorn の反復は以下のように読める。
\begin{align}
	\label{eq-lse-dual-1}\itt{\fD} &= \SMINEr\, (\C-\ones_n\transp{\it{\gD}}) + \varepsilon \log \a, \\
	\label{eq-lse-dual-2}\itt{\gD} &= \SMINEc\, (\C-\it{\fD}\transp{\ones_m})  + \varepsilon \log \b.
\end{align}
$\varepsilon \rightarrow 0$ のとき $\smine$ は $\min$ に収束するが、$\varepsilon=0$ の極限では反復はもはや収束しない。なぜなら、制約付き問題の場合には交互最小化は収束せず、正則化なし双対~\eqref{eq-dual}がそれに該当するからである。
\end{rem}
%%%%%%%%%%



%%%%%%%%%%
\begin{rem}[対数領域 Sinkhorn]\label{rem-log-sinkh}
Sinkhorn の更新~\eqref{eq-sinkhorn}と数学的に等価ではあるが、反復~\eqref{eq-slse-smin-1}および~\eqref{eq-slse-smin-2}は、$\varepsilon$ の値が小さい場合のアンダーフローを回避するための\emph{対数-和-指数}安定化トリックの使用を示唆する。$\underbar{z}=\min \z$ と書くと、そのトリックは $\smine \z$ を以下のように評価することを示唆する。
\eql{\label{eq-log-domain-min}
	\smine \z = \underbar{z} -\varepsilon \log \sum_i e^{-(\z_i-\underbar{z})/\varepsilon}.
}
~\eqref{eq-log-domain-min}のように $\underbar{z}$ を差し引いて対数領域反復を安定化させる代わりに、実際には以前に計算されたスケーリングを差し引くことができる。
これにより安定化された反復が得られる。
\begin{align}
	\label{eq-lse-sinkh-1}\itt{\fD} &= \SMINEr( \logP(\it{\fD},\it{\gD}))  + \it{\fD} + \varepsilon\log(\a), \\
	\label{eq-lse-sinkh-2}\itt{\gD} &= \SMINEc( \logP(\itt{\fD},\it{\gD})) + \it{\gD} + \varepsilon \log(\b),
\end{align}
ここで以下を定義した。
\eq{
	\logP(\fD,\gD) = \pa{\C_{i,j} - \fD_i - \gD_j}_{i,j}.
}
もとの反復~\eqref{eq-sinkhorn}とは対照的に、これらの対数領域反復~\eqref{eq-lse-sinkh-1}および~\eqref{eq-lse-sinkh-2}は任意の $\varepsilon>0$ に対して安定である。
なぜなら量 $\logP(\fD,\gD)$ は反復中に有界にとどまるからである。
欠点は、各ステップで $nm$ 回の $\exp$ の計算を必要とすることである。
$\SMINEr$ や $\SMINEc$ の計算は、行列乗算よりも典型的にかなり遅く、行列 $\logP$ の行ごとのソフト最小値を計算する必要がある。したがって、複数の周辺分布に対する Sinkhorn 写像の適用を効率的に並列化する方法は存在しない。
小次元のユークリッド領域では、$\varepsilon$ を減衰させる戦略とスパースグリッドを組み合わせた効率的なマルチスケールソルバーを開発することで、計算を大幅に高速化することが可能である~\citep{schmitzer2016stabilized}。
\end{rem}
%%%%%%%%%%




%%%%%%%%%%
\begin{rem2}{一般的な測度に対する双対}
一般的な（必ずしも離散でない）入力測度 $(\al,\be)$ に対して、双対問題~\eqref{eq-dual-formulation}は以下のように読める。
\eql{\label{eq-dual-entropic}
	\usup{(\f,\g) \in \Cc(\X)\times\Cc(\Y)} \int_\X \f \d\al + \int_\Y \g \d\be
		 - \varepsilon \int_{\X\times\Y} e^{ \frac{-c(x,y)+f(x)+g(y)}{\varepsilon} } \d\al(x)\d\be(y).
}
これはもとの問題~\eqref{eq-dual-generic}に現れる制約 $\Potentials(\c)$ の平滑化に対応し、$\varepsilon \rightarrow 0$ の極限で復元される。
エントロピー輸送の場合にこれらの Kantorovich ポテンシャル $(\f,\g)$ の存在を証明する（\ie sup が実際に max である）ことは、古典的 OT の場合より容易ではない。なぜなら、$c$-変換を用いることができず、ポテンシャルは自動的には Lipschitz にならないからである。存在の証明は Sinkhorn 反復の収束を用いて行うことができる。詳細については~\citep{2016-chizat-sinkhorn}を参照されたい。
\end{rem2}
%%%%%%%%%%



%%%%%%%
\begin{rem2}{制約なしエントロピー双対}
注意~\ref{rem-uncons-dual}と同様に、$\int_\X \d\mu=\int_\Y\d\nu = 1$ の場合、以下の代替的な双対定式化を考えることができる。
\eql{\label{eq-dual-entropic-alt}
	\usup{(\f,\g) \in \Cc(\X)\times\Cc(\Y)} \int_\X \f \d\al + \int_\Y \g \d\be
		+ \smine( c - f \oplus g ),
}
これは~\eqref{eq-dual-entropic}と同じ最適値を達成する。
~\eqref{eq-defn-softmin}と同様に、ソフト最小値（ここでは $\X \times \Y$ 上の）は以下で定義される。
\eq{
	\foralls S \in \Cc(\X\times\Y), \quad
	\smine S \eqdef
	- \varepsilon \int_{\X\times\Y} e^{ \frac{-S(x,y)}{\varepsilon} } \d\al(x)\d\be(y)
}
（これは $(\al,\be)$ に依存することに注意されたい）。
$\epsilon \rightarrow 0$ のとき、$\smine \rightarrow \min$ であり、正則化なしの制約なし定式化~\eqref{eq-dual-uncons}で用いられるものに帰着する。
~\eqref{eq-dual-entropic}と~\eqref{eq-dual-entropic-alt}のいずれもが制約なし問題であるが、~\eqref{eq-dual-entropic-alt}の主な利点はより良い条件付けを持つことである。汎関数の Hesse 行列が $\epsilon$ で一様に有界であるという意味においてである。
このような条件付けの改善を得る別の方法は、半双対問題を考えることである。第~\ref{sec-semi-discr-entropy}節、特に注意~\ref{rem-second-order-smoothness}を参照されたい。
この代替的な双対定式化の欠点は、対数の存在が第~\ref{sec-sgd}節で詳述される確率的最適化法の使用を妨げることである。特に注意~\ref{rem-sgd-cont-cont}を参照されたい。
\end{rem2}
%%%%%%%

%%%%%%%%%%%%%%%%%%%%%%%%%%%%%%%%%%%%%%%%%%%%%%%%%%%%%%%%%%%%%%%%%%%%%%%%%%
\section{Optimal Transport コストの正則化近似}
\label{sec-regularized-cost}

エントロピー双対~\eqref{eq-dual-formulation}は滑らかな制約なし凹最大化問題であり、もとの Kantorovich 双対~\eqref{eq-dual}を近似する。これは以下の命題で詳述される。

\begin{prop}\label{prop-feasibility-dual}
~\eqref{eq-dual-formulation}の任意の最適解の対 $(\fD^\star,\gD^\star)$ は $(\fD^\star,\gD^\star) \in \PotentialsD(\C)$（~\eqref{eq-feasible-potential}で定義される実行可能 Kantorovich ポテンシャルの集合）を満たす。結果として、任意の $\varepsilon$ に対して
\eq{
	\dotp{\fD^\star}{\a} + \dotp{\gD^\star}{\b} \leq  \MKD_\C(\a,\b).
}
\end{prop}
\begin{proof}
	\eqref{eq-prim-dual-ent}における主双対最適性条件と $\P$ が確率であるため $\P_{i,j}\leq 1$（すべての $i,j$ に対して）という制約から、$\exp(-(\fD^\star_i+\gD^\star_j-\C_{i,j})/\varepsilon)\leq 1$ が得られ、したがって $\fD^\star_i+\gD^\star_j\leq\C_{i,j}$ が成り立つ。
\end{proof}

~\eqref{eq-regularized-discr}で定義される正則化輸送コスト $\MKD_\C^\varepsilon$ の重要な利点は、滑らかかつ凸であることであり、変分問題における損失関数としての統合に最適である（第~\ref{c-variational}章を参照）。

\begin{prop}\label{prop-convexity-dual}
$\MKD_\C^\varepsilon(\a,\b)$ は $\epsilon\geq 0$ に対して $\a$ と $\b$ の同時凸関数である。$\epsilon>0$ のとき、その勾配は以下に等しい。
\eq{
	\nabla\MKD_\C^\varepsilon(\a,\b)= \begin{bmatrix} \fD^\star \\ \gD^\star \end{bmatrix},
}
ここで $\fD^\star$ と $\gD^\star$ は座標の和が $0$ になるように選ばれた式~\eqref{eq-dual-formulation}の最適解である。
\end{prop}


\citep{CuturiSinkhorn}において、2つのヒストグラム間の Wasserstein 距離を近似するための上下界が提案された。これらの界は、Sinkhorn アルゴリズムが提供する解における主問題と双対問題の目的関数の値の評価からなる。

\begin{defn}[Sinkhorn ダイバージェンス]
	$\fD^\star$ と $\gD^\star$ を~\eqref{eq-dual-formulation}の最適解、$\P^\star$ を~\eqref{eq-regularized-discr}の解とする。Wasserstein 距離は以下の主および双対 Sinkhorn ダイバージェンスを用いて近似される：
\begin{align*}
	\SINKHORNP_\C^\varepsilon(\a,\b)& \eqdef \dotp{\C}{\P^\star} =  \dotp{e^{\frac{\fD^\star}{\varepsilon}}}{(\K\odot \C)e^{\frac{\gD^\star}{\varepsilon}}}, \\
	\SINKHORND_\C^\varepsilon(\a,\b) & \eqdef \dotp{\fD^\star}{\a} + \dotp{\gD^\star}{\b},
\end{align*}
ここで $\odot$ は行列の要素ごとの積を表す。
\end{defn}

\begin{prop}\label{prop-sinkhorn-div}
以下の関係が成り立つ：
\eq{
	\SINKHORND_\C^\varepsilon(\a,\b)\leq \MKD_\C^\varepsilon(\a,\b) \leq \SINKHORNP_\C^\varepsilon(\a,\b).
}
さらに
\eql{\label{eq-sinkh-div-gap}
	\SINKHORNP_\C^\varepsilon(\a,\b) - \SINKHORND_\C^\varepsilon(\a,\b) = \varepsilon ( \HD(\P^\star)+1 ) .
}
\end{prop}
\begin{proof}
式~\eqref{eq-sinkh-div-gap}は、最適解において主問題と双対問題が同じ値を持つことを書くことで得られる（~\eqref{eq-dual-formulation}を参照）。したがって
\eq{
	\MKD_\C^\varepsilon(\a,\b) = \SINKHORNP_\C^\varepsilon(\a,\b) - \varepsilon \HD(\P^\star) =
	\SINKHORND_\C^\varepsilon(\a,\b) - \varepsilon \dotp{e^{\fD^\star/\varepsilon} }{ \K e^{\gD^\star/\varepsilon}}
}
最後の結果は $\dotp{e^{\fD^\star/\varepsilon} }{ \K e^{\gD^\star/\varepsilon}}=1$ であることに気づくことで得られる。後者は $\P^\star$ のすべての成分の和を計算することに等しいからである。
\end{proof}


上で与えた関係は実際の OT 距離を近似する実用的な方法を示唆するが、実は Sinkhorn アルゴリズムの収束時にのみ有効であり、したがって実際には本当に有用とはいえない。実際には、Sinkhorn 反復は常にある精度閾値に到達した後に打ち切られる。所定の $L$ 回の反復を設定し、最適解 $\fD^\star$ と $\gD^\star$ の代わりに反復 $\itL{\fD}$ と $\itL{\gD}$ を用いて $\SINKHORND_\C^\varepsilon$ を評価する場合、下界が復元される：式~\eqref{eq-lse-sinkh-1}および~\eqref{eq-lse-sinkh-2}に現れる記法を用いて、$\MKD_\C^\varepsilon$ の以下の有限ステップ近似を導入する：
\eql{\label{eq-algorithmic-loss}
	\itL{\SINKHORND_\C}(\a,\b) \eqdef \dotp{\itL{\fD}}{\a} + \dotp{\itL{\gD}}{\b}.
}
この「アルゴリズム的」Sinkhorn 汎関数は $L\geq 1$ であれば正則化コスト関数の下界を与える。


\begin{prop}[有限 Sinkhorn ダイバージェンス]\label{defn-sinkhorn-div}
以下の関係が成り立つ：
\eq{
	\itL{\SINKHORND_\C}(\a,\b) \leq \MKD_\C^\varepsilon(\a,\b).
}
\end{prop}
\begin{proof}命題~\ref{prop-feasibility-dual}の証明と同様に、たった1回の Sinkhorn 反復の後でさえ、~\eqref{eq-slse-sinkh-1}および~\eqref{eq-slse-sinkh-2}に従って $\itL{\fD}$ と $\itL{\gD}$ は、成分 $\exp(-(\itL{\fD}_i+\itL{\gD}_j-\C_{i,j})/\varepsilon)$ からなる行列の列和が $\b$ であり、したがってその成分はすべて $1$ で上から抑えられるという事実を利用する。これにより $(\itL{\fD}_i,\itL{\gD})$ の双対実行可能性が得られる。
\end{proof}

\begin{rem}[Sinkhorn 反復の主実行不可能性]
\eqref{eq-sink-matrix}で与えられる主反復は主実行可能ではないことに注意されたい。定義上、これらの反復は収束時に周辺制約を満たすように設計されているためである。したがって、$\P^{(2L+1)}$ は実行可能でないため、$\dotp{\C}{\P^{(2L+1)}}$ を $\MKD_\C(\a,\b)$ の近似として考えることは妥当でない。注意~\ref{rem-complexity-rounding}で示された \citet{altschuler2017near} の丸めスキームを用いることで、$\MKD_\C^\varepsilon(\a,\b)$ の上界が得られ、さらにこれは Sinkhorn アルゴリズムと同様に複数のヒストグラムのペアに対して行列演算を用いて並列に便利に計算できる~\citep{NIPS2018_8184}。
\end{rem}


\begin{rem}[有限双対 Sinkhorn ダイバージェンスの非凸性] ~\eqref{eq-dual-formulation}の正則化された式 $\MKD_\C^\varepsilon$ とは異なり、有限 Sinkhorn ダイバージェンス $\itL{\SINKHORND_\C}(\a,\b)$ は一般にその引数の凸関数では\emph{ない}（これは数値的に容易に確認できる）。しかし $\itL{\SINKHORND_\C}(\a,\b)$ は微分可能な関数であり、自動微分技法（注意~\ref{rem-auto-diff}を参照）を用いてその任意の引数、特に $\C,\a$、または $\b$ に関して微分できる。
\end{rem}




%%%%%%%%%%%%%%%%%%%%%%%%%%%%%%%%%%%%%%%%%%%
\section{一般化 Sinkhorn}
\label{sec-generalized}

正則化 OT 問題~\eqref{eq-regularized-discr}は、以下の形の構造化凸最適化問題の特殊ケースである。
\eql{\label{eq-generalized-ot-regul}
	\umin{\P}
		\sum_{i,j} \C_{i,j} \P_{i,j} - \varepsilon \HD(\P) + F( \P\ones_m ) + G( \transp{\P}\ones_n ).
}
実際、$F=\iota_{\{\a\}}$ および $G=\iota_{\{\b\}}$ と定義すると、ここで閉凸集合 $\Cc$ の指示関数は
\eql{\label{eq-iota-function}
	\iota_{\Cc}(x) = \choice{
		0 \qifq x \in \Cc, \\
		+\infty \quad\text{otherwise},
	}
}
であり、$\Couplings(\a,\b)$ を定義するハードな周辺制約が復元される。
命題~\ref{prop-regularized-primal}の証明はこのより一般的な問題~\eqref{eq-generalized-ot-regul}にも引き継がれ、~\eqref{eq-generalized-ot-regul}の一意の解もまた~\eqref{eq-scaling-form}の形を持つ。

\citep{2015-Peyre-siims,FrognerNIPS,2016-chizat-sinkhorn,karlsson2016generalized}で示されるように、Sinkhorn 反復~\eqref{eq-sinkhorn}はこの問題に拡張でき、以下のようになる。
\eql{\label{eq-gen-sinkh}
	\uD \leftarrow \frac{\Prox_{F}^{\KLD}(\K \vD)}{\K \vD}
	\qandq
	\vD \leftarrow \frac{\Prox_{G}^{\KLD}(\transp{\K}\uD)}{\transp{\K}\uD},
}
ここで $\KLD$ ダイバージェンスに対する近接作用素は
\eql{\label{eq-prox-kl}
	\foralls \uD \in \RR_+^N, \quad
	\Prox_{F}^{\KLD}(\uD) = \uargmin{\uD' \in \RR_+^N}	\KLD(\uD'|\uD) + F(\uD').
}
いくつかの関数 $F,G$ に対して反復~\eqref{eq-gen-sinkh}の線形収束率を証明することが可能であり、これらのスキームは任意の測度に一般化できる。詳細は~\citep{2016-chizat-sinkhorn}を参照されたい。

反復~\eqref{eq-gen-sinkh}が興味深いのは、$\Prox_{F}^{\KLD}$ と $\Prox_{G}^{\KLD}$ が閉じた形で、または非常に効率的に計算できる場合である。これは特に $F(\uD) = \sum_i F_i(\uD_i)$ の形の分離可能関数の場合に当てはまる。この場合
\eq{
	\Prox_{F}^{\KLD}(\uD) = \pa{ \Prox_{F_i}^{\KLD}(\uD_i) }_i.
}
各 $\Prox_{F_i}^{\KLD}$ の計算は通常単純である。スカラー最適化問題であるためである。
もとの Sinkhorn アルゴリズムと同様に、対数領域での計算を用いて計算を安定化することも可能である~\citep{2016-chizat-sinkhorn}。

このアルゴリズムは OT の様々な一般化の解を近似するために用いることができ、特に~\eqref{eq-unbalanced-pbm}の形の非平衡 OT 問題（第~\ref{sec-unbalanced}節、特に反復~\eqref{eq-iterate-gen-sinkh}を参照）および~\eqref{eq-grad-flow-discr}の形の勾配流問題（第~\ref{sec-grad-flows}節を参照）に用いることができる。


%%%
\begin{rem1}{双対性と Legendre 変換}
~\eqref{eq-generalized-ot-regul}の双対問題は以下のようになる。
\eql{\label{eq-dual-generalized}
	\umax{\fD,\gD}
		- F^*( \fD ) -G^*(\gD)
		- \varepsilon\sum_{i,j} e^{ \frac{\fD_i+\gD_j-\C_{i,j}}{\varepsilon} }
}
したがって $(\uD,\vD) = (e^{\fD/\varepsilon},e^{\gD/\varepsilon})$ が~\eqref{eq-scaling-form}に現れる関連するスケーリングである。
ここで $F^*$ と $G^*$ は Fenchel--Legendre 共役であり、以下で定義される凸関数である。
\eql{\label{eq-legendre}
		\foralls \fD \in \RR^n, \quad
		F^*(\fD) \eqdef \umax{\a \in \RR^n} \dotp{\fD}{\a} - F(\a).
}
一般化 Sinkhorn 反復~\eqref{eq-gen-sinkh}は Dykstra のアルゴリズム~\citep{Dykstra83,Dykstra85}（Bregman ダイバージェンスに拡張されたもの~\citep{bauschke-lewis,CensorReich-Dykstra}。注意~\ref{rem-bregman}も参照）の特殊ケースであり、双対問題~\eqref{eq-dual-generalized}に対する交互最大化スキームである。
\end{rem1}
%%%

定式化~\eqref{eq-generalized-ot-regul}は、2つより多くの関数と単一のカップリングを超えてさらに一般化できる。詳細は~\citep{2016-chizat-sinkhorn}を参照されたい。これには特殊ケースとして、第~\ref{sec-multimarginal}節で詳述される多周辺問題に対する Sinkhorn アルゴリズム~\eqref{eq-sinkh-multimarg}が含まれる。
正則化重心問題~\eqref{eq-entropic-bary}もこの方法で書き直すことが可能であり、反復~\eqref{eq-sinkhorn-bary}は実際にこの一般化 Sinkhorn の特殊ケースである。


\if 0 %% REMOVED %%%
\eq{
	\umin{(\P_s)_{s=1}^S}
		\sum_{i,j,s} \C_{s,i,j} \P_{s,i,j} - \varepsilon \HD(\P_s) +
		\sum_{k} F_k( (\P_s\ones_m)_s, (\transp{\P_s}\ones_n)_s )
}
where $\{\C_s\}_s$ are given cost matrices.

The solution has the form
\eql{\label{eq-scaling-form-multi}
	\foralls (s,i,j) \in \range{1,S} \times \range{n} \times \range{m}, \quad
	\P_{s,i,j} =  \K_{s,i,j} \prod_{k} \uD_{k,i} \prod_\ell \vD_{k,j}
}
where $\K_s=e^{-\C_s/\varepsilon}$ is a Gibbs kernel, and $(\uD_k,\vD_\ell)_{k,\ell}$ are scaling vectors that need to be computed using Generalized Sinkhorn iterations similar to~\eqref{eq-gen-sinkh}~\citep{2016-chizat-sinkhorn}.
%
In the case where the functions $F_s$ are indicator functions, one retrieves the Sinkhorn algorithm~\eqref{eq-sinkh-multimarg} for the multi-marginal problem, as detailed in~\S\ref{sec-multimarginal}.
%
It is also possible to rewrite the regularized barycenter problem~\eqref{eq-entropic-bary} this way, and the iterations~\eqref{eq-sinkhorn-bary} are in fact a special case of this generalized Sinkhorn.

\fi %%%%
